Put \(q_0=\min\{\beta,\gamma\}\). We first record why both arm regressions have the smoothness needed below. For a member \(P\in\mathcal M_{\mathrm{tr}}(\gamma;\kappa)\), calibration and the two margin inequalities at \(t=t_0=c\) give
\[
 P_X\{0<|\tau_P(X)-\theta|\le c\}=1.
\]
The density is bounded below on the whole cube and \(\tau_P\) is continuous. Hence \(|\tau_P(x)-\theta|\le c\) for every \(x\in\mathcal X\), the zero level has \(P_X\)-measure zero, and \(\mathcal D_P(c)=\mathcal X\). Thus \(\mathrm{ass:near-band-holder}\) is a global \(\gamma\)-Hölder condition. By \(\mathrm{def:cate}\), the continuous version of the treated regression is \(\mu_1=\mu_0+\tau_P\). The elementary embedding of Hölder balls therefore shows that both \(\mu_0\) and \(\mu_1\) belong to a \(q_0\)-Hölder ball whose radius is bounded uniformly over the calibrated class.

We now define the estimator. Let \(p_0=\lceil q_0\rceil-1\), use the window, kernel, and shifted-Legendre vector from \(\mathrm{def:window-basis}\) with order \(p_0\), and write \(r(u)=r_{x,h,p_0}(u)\). For \(a\in\{0,1\}\), set
\[
 \widetilde Q_a(x)=\frac1n\sum_{i=1}^n K_{x,h}(X_i)r(X_i)r(X_i)^\top\mathbf 1\{A_i=a\},
 \qquad
 \widetilde R_a(x)=\frac1n\sum_{i=1}^n K_{x,h}(X_i)r(X_i)Y_i\mathbf 1\{A_i=a\}.
\]
Let \(e_1(x)=\pi(x)\) and \(e_0(x)=1-\pi(x)\). Under \(\mathrm{lem:uniform-design}\), the population matrix is
\[
 Q_a(x)=\mathbb E_P\!\left[K_{x,h}(X)r(X)r(X)^\top\mathbf 1\{A=a\}\right].
\]
For every vector \(v\), the change of variables \(z=T_{x,h}(u)\), orthonormality of the Legendre basis, and strict overlap give
\[
 2^{-d}c_\pi\|v\|_2^2
 \le v^\top Q_a(x)v
 \le 2^{-d}(1-c_\pi)\|v\|_2^2.
\]
Clip the eigenvalues of the symmetric matrix \(\widetilde Q_a(x)\) to the fixed interval \([2^{-d-1}c_\pi,2^{1-d}(1-c_\pi)]\), call the result \(\widehat Q_a(x)\), and define
\[
 \widehat\mu_a(x)=\operatorname{clip}_{[-1,1]}
 \left\{r_x^\top\widehat Q_a(x)^{-1}\widetilde R_a(x)\right\},
 \qquad
 \widehat\tau^{\mathrm{dir}}_n(x)=\widehat\mu_1(x)-\widehat\mu_0(x),
\]
\[
 \widehat G_n^{\mathrm{dir}}=\{x\in\mathcal X:\widehat\tau^{\mathrm{dir}}_n(x)>\theta\}.
\]
This is a single measurable map of the sample; spectral clipping makes every inverse well defined.

We next prove its pointwise exponential bound. Let \(t_{a,x}\) be the degree-\(p_0\) Taylor polynomial of \(\mu_a\) on \(W_{x,h}\), expressed in the vector \(r\). The Hölder remainder and \(h/2\le\ell_j\le h\) imply
\[
 \sup_{u\in W_{x,h}}|\mu_a(u)-t_{a,x}(u)|\le C h^{q_0}.
\]
Consequently, if \(b_{a,x}\) is the coefficient vector of \(t_{a,x}\), then
\[
 \left\|\mathbb E_P\widetilde R_a(x)-Q_a(x)b_{a,x}\right\|_2\le C h^{q_0},
 \qquad
 |r_x^\top b_{a,x}-\mu_a(x)|=0.
\]
All coordinates of \(r\) are bounded because their dimension depends only on \(d\) and \(q_0\). Each centered entry of \(K_{x,h}rr^\top\mathbf 1\{A=a\}\) and of \(K_{x,h}rY\mathbf 1\{A=a\}\) has variance at most \(Ch^{-d}\) and envelope at most \(Ch^{-d}\). Indeed, \(K_{x,h}\le C h^{-d}\), its support has probability at most \(Ch^d\), and \(|Y|\le1\). For such a centered variable \(S\),
\[
 \mathbb E|S|^m\le (Ch^{-d})^{m-2}\mathbb E S^2\qquad(m\ge2).
\]
Expanding the exponential series, summing the resulting geometric series for \(|\lambda|Ch^{-d}<1\), and optimizing the exponential Markov inequality yields Bernstein's bound
\[
 P\!\left(\left|\frac1n\sum_{i=1}^nS_i\right|>
 C\left\{\sqrt{z/(nh^d)}+z/(nh^d)\right\}\right)\le2e^{-z}.
\]
A union bound is harmless because the polynomial dimension is fixed.

Eigenvalue clipping is the orthogonal Frobenius projection onto the closed convex spectral interval. If \(S=\widetilde Q_a(x)\), \(U=\widehat Q_a(x)\), and \(Q=Q_a(x)\), its variational inequality gives \(\langle S-U,Q-U\rangle_F\le0\). Since \(Q\) lies in the interval,
\[
 \|U-Q\|_F^2
 =\langle U-S,U-Q\rangle_F+\langle S-Q,U-Q\rangle_F
 \le\|S-Q\|_F\|U-Q\|_F.
\]
Thus \(\|\widehat Q_a-Q_a\|_F\le\|\widetilde Q_a-Q_a\|_F\), and every clipped inverse is uniformly bounded. Using
\[
 r_x^\top\widehat Q_a^{-1}\widetilde R_a-r_x^\top b_{a,x}
 =r_x^\top\widehat Q_a^{-1}
 \{\widetilde R_a-\widehat Q_a b_{a,x}\}
\]
and the preceding bias and concentration bounds gives, uniformly in \(P\), \(x\), and \(a\),
\[
 P\!\left[|\widehat\mu_a(x)-\mu_a(x)|>
 C\left\{h^{q_0}+\sqrt{z/(nh^d)}+z/(nh^d)\right\}\right]
 \le Ce^{-cz}.
\]
Final scalar clipping cannot enlarge the error because \(|\mu_a|\le1\).

Choose \(h=n^{-1/(2q_0+d)}\) and put \(\varepsilon_n=h^{q_0}=n^{-q_0/(2q_0+d)}\). Then \((nh^d)^{-1/2}=\varepsilon_n\). For \(1\le z\le c/\varepsilon_n\), the preceding display and a union bound over the two arms imply
\[
 P\{|\widehat\tau^{\mathrm{dir}}_n(x)-\tau_P(x)|>Cz\varepsilon_n\}
 \le Ce^{-cz}.
\]

It remains to transfer this pointwise inequality to the weighted set loss. Write \(\Delta(x)=|\tau_P(x)-\theta|\). Misclassification implies \(|\widehat\tau^{\mathrm{dir}}_n(x)-\tau_P(x)|\ge\Delta(x)\). For a sufficiently large fixed \(A\), the part with \(0<\Delta(X)\le A\varepsilon_n\) is at most
\[
 A\varepsilon_n P_X\{0<\Delta(X)\le A\varepsilon_n\}
 \le C\varepsilon_n^2
\]
by \(\mathrm{ass:upper-margin}\) with exponent one. On the dyadic shell
\[
 2^jA\varepsilon_n<\Delta(X)\le t_j,
 \qquad
 t_j=\min\{2^{j+1}A\varepsilon_n,c\},
\]
the pointwise tail bound and the margin inequality at \(t_j\) give a contribution at most
\[
 Ct_j^2e^{-c2^j}
 \le C(2^j\varepsilon_n)^2e^{-c2^j}.
\]
The series \(\sum_{j\ge0}4^je^{-c2^j}\) is finite, and \(\Delta\le c\), so these truncated shells exhaust the support. Tonelli's theorem now yields
\[
 \sup_{P\in\mathcal M_{\mathrm{tr}}(\gamma;\kappa)}
 \mathbb E_P[d_H(\widehat G_n^{\mathrm{dir}},G_P;P)]
 \le C\varepsilon_n^2.
\]
Taking the infimum over all measurable set estimators in \(\mathrm{def:minimax-risk}\) proves the asserted upper bound.